\subsection{\resultrevision{Shared and Task-Specific Tree Structures}}
\label{sec:learned-tree-structures}
Figure~\ref{fig:learned-sequencing-trees} shows the sequencing trees learned together in R4 for scenarios~1, 3 and~5, which share the maximum-flow-time and maximum-weighted-tardiness objectives at utilisations 0.70, 0.80 and 0.90. The case uses representative seed~20 and the first training-ranked generation-49 heuristic for each task. The trees contain 37, 31 and 71 nodes.

Two modules recur across tasks with identical operators, terminals and ordered child structures. A, $100\,\mathrm{NOR}\,w_2-\mathrm{SL}$, is a five-node subtree shared by S1 and S3; B, $\mathrm{W}+100w_2$, is a three-node subtree shared by S3 and S5. The preference terminals supply $100w_1$ and $100w_2$, where $w_1+w_2=1$. A combines remaining-operation count, preference and slack, whereas B combines job weight and preference. Sharing is pair-specific: no nonterminal subtree occurs in all three trees.

Each task also retains distinctive combinations. U1, $\max(\mathrm{TIS},\mathrm{WR}-\mathrm{TR})$, links time in system with transfer-adjusted remaining work in S1. U3, $\min(\mathrm{NIQ},\mathrm{SL})$, appears three times in S3, coupling queue length and slack in its competing priority branches. U5, $D(\mathrm{WR},100w_1)+\mathrm{NIQ}+\mathrm{MWT}$, combines preference-scaled remaining work, queue length and machine waiting time in S5; $D$ denotes protected division. These modules are unique to their respective trees in this case. The trees contain 9, 9 and 31 distinct proper nonterminal subtrees absent from the other two trees. S1 and S5 use subtractive roots, while S3 selects between branches through a minimum. Shared fragments therefore coexist with different complete decision rules.

This pattern reflects the form of reuse targeted by MPSLGP: donor fragments are assessed within receiving-task contexts, while task-local evaluation and balanced quotas retain separate heuristics. Node contributions assign each child's task, and offspring outcomes update the corresponding task/region utility. In this run, S1-to-S3 transfer produced 5,039 tracked children and 1,167 survivors; the reverse direction produced 4,902 and 1,278. S3-to-S5 and S5-to-S3 produced 5,536/1,137 and 5,711/1,328 children/survivors, respectively. These generation-15--49 records show that both receiving directions retained cross-task offspring. The available records establish exchange activity and structural overlap, but do not resolve the genealogy of individual modules. The case thus illustrates reusable preference-aware structure alongside task-specific adaptation.

\begin{figure*}[t]
\centering
\includegraphics[width=\textwidth]{results/latest-20261106/tree-case-study/learned-sequencing-trees.pdf}
\caption{\resultrevision{Generation-49 sequencing trees for R4, run~20: S1 and S3 share the upper row; S5 spans the lower row. All 139 nodes and their child order are retained. Circles denote operators and boxes denote terminals. Light blue A and light orange B mark shared subtrees; neutral grey U modules mark task-specific examples. Preference inputs $w_1$ and $w_2$ use blue and orange, respectively. Their terminal values are $100w_1,100w_2$. WR = WKR and TR = TRANT.}}
\label{fig:learned-sequencing-trees}
\end{figure*}

\subsection{\resultrevision{Task-Specific Effects of Preference Inputs}}
\label{sec:tree-preference-effects}
The three trees assign different roles to the same preference inputs. The weights $w_1$ and $w_2=1-w_1$ express preferences for maximum flow time and maximum weighted tardiness, respectively, and enter the trees through terminal values $100w_1$ and $100w_2$. Scheduling attributes retain their raw values, and smaller priority scores are selected first. S1 and S3 each use two $w_2$ leaves; S5 uses eight $w_2$ and two $w_1$ leaves. Because the weights are complementary, $w_2$ alone can condition a heuristic. Algebraic reduction reveals how these inputs interact with task-state features and which contributions affect ranking.

For S1, write $X=\mathrm{WR}-\mathrm{TR}$. The repeated $X-X$ branches cancel, giving
\begin{equation}
\begin{aligned}
g_1={}&\min(\mathrm{SL}-100\,\mathrm{NOR}\,w_2,-X)\\
&-\max(\mathrm{TIS},X)-\min(100w_2,\mathrm{SL}).
\end{aligned}
\label{eq:case-tree-s1}
\end{equation}
Preference interacts with remaining-operation count and slack. Increasing $w_2$ changes the first minimum's active branch according to the operation's slack and remaining work. The final minimum introduces an additional slack-dependent threshold. These interactions give different operations different responses to the same preference. The cancelled branches also reveal redundant syntax in the evolved expression.

For S3, let $Q=\min(\mathrm{NIQ},\mathrm{SL})$. Its tree reduces to
\begin{equation}
\begin{aligned}
g_3&=\min\{C_3+100(1-\mathrm{NOR})w_2,\ H_3\},\\
C_3&=Q-\mathrm{WR}+\mathrm{W}+\mathrm{SL},\\
H_3&=\mathrm{W}[Q-\mathrm{WR}+\min(\mathrm{PT},\mathrm{WR})]-Q.
\end{aligned}
\label{eq:case-tree-s3}
\end{equation}
A and B meet in the first branch, where their preference terms combine into $100(1-\mathrm{NOR})w_2$. This contribution vanishes for $\mathrm{NOR}=1$ and lowers the branch score as the tardiness weight increases when $\mathrm{NOR}>1$. The outer minimum can instead select the preference-free branch $H_3$. The same A fragment therefore has a different effective role in S3 than in S1 because its surrounding expression differs, illustrating the context-dependent reuse enabled by subtree crossover.

S5 couples both preferences with waiting information. Define $D(a,b)=1$ for $b=0$, and $D(a,b)=a/b$ otherwise. Within one work-centre queue, NIQ and the supplied weights are shared across candidates. Terms involving only these quantities are common offsets. Removing them gives the ranking-equivalent form
\begin{equation}
\begin{aligned}
g_5&\overset{\mathrm{rank}}{\sim}
C_5-100\,\mathrm{W}w_2\\
&\quad-[D(\mathrm{WR},100w_1)+\mathrm{NIQ}+\mathrm{MWT}]K_5,\\
K_5&=D\!\left(\max(\mathrm{NOR},100w_2),\mathrm{NOR}\,\mathrm{OWT}\right),
\end{aligned}
\label{eq:case-tree-s5}
\end{equation}
where $C_5$ contains the remaining preference-free terms. The $-100\,\mathrm{W}w_2$ term favours higher-weight jobs as the tardiness weight increases, all else equal. U5 and $K_5$ couple preference with remaining work and queue and waiting conditions. The scaled $w_1$ input enters a denominator; at $w_1=0$, protected division returns 1 for that subexpression. B is multiplied by $100w_2$ in S5, in contrast to its additive placement in S3. Combining the complete expression retains the job-weight interaction while eliminating several preference-only offsets. The task-specific context determines which preference contributions remain active in ranking.

Each fixed dual-tree heuristic produces 200 distinct objective pairs across the 200 test preferences. S1 and S5 receive explicit preferences through sequencing, with routing responding indirectly to the resulting shop states; S3 also uses $w_1$ in routing. The system-level responses therefore reflect both rules and their interaction. The region-aware controller acts during evolution, while preference terminals condition the deployed decisions. Together, they support cross-task exchange of preference-aware fragments and task-specific responses through different surrounding structures.